\section{Caracterización dinámica y parámetros}

El comportamiento libre está gobernado por la ecuación característica
\begin{equation}
  ms^2+\beta s+k=0,
\end{equation}
cuyas raíces son
\begin{equation}
  s_{1,2}=\frac{-\beta\pm\sqrt{\beta^2-4mk}}{2m}.
  \label{eq:polos}
\end{equation}
Según el valor de $\beta^2-4mk$, el sistema será subamortiguado, críticamente amortiguado o sobreamortiguado. En términos de $\zeta$, estos regímenes corresponden a $0<\zeta<1$, $\zeta=1$ y $\zeta>1$, respectivamente. Si $m>0$, $\beta>0$ y $k>0$, los polos presentan parte real negativa y el equilibrio es asintóticamente estable.

Para el caso subamortiguado se define la frecuencia natural amortiguada
\begin{equation}
  \omega_d=\omega_n\sqrt{1-\zeta^2}.
\end{equation}
La respuesta natural tendrá entonces la forma general
\begin{equation}
  y_n(t)=e^{-\zeta\omega_nt}
  \left[A\cos(\omega_dt)+B\sin(\omega_dt)\right],
  \label{eq:respuesta_natural}
\end{equation}
donde $A$ y $B$ se determinan a partir de las condiciones iniciales.

\draftnote{Seleccionar y justificar los valores de $m$, $k$ y $\beta$; calcular $\omega_n$, $\zeta$, $\omega_d$ y los polos; completar una tabla con símbolo, valor, unidad, fuente y criterio físico.}

